% GENERATED FILE. Edit canonical lessons, then run npm run generate:books.
% canonical-source-hash: fnv1a64:aeb52a5738d398da
% canonical-lessons: FA-C428-aber-e-piyade, FA-C428-motorsiklet, FA-C428-kamiyun, FA-C428-vanet, FA-C428-zir-gereftan

\chapter{Other road users}
\label{ch:fa-other-road-users}
\hlchaptermodality{428}

\begin{chapteropening}
\textbf{By the end of this chapter:} \emph{I can talk about other road users: pedestrians, motorbikes, lorries, pickups and running someone over.}

Complete the last lesson of chapter 428: I can talk about other road users: pedestrians, motorbikes, lorries, pickups and running someone over.
\end{chapteropening}

\section[\texorpdfstring{\fa{عابر} \fa{پیاده} ('âber-e piyâde)}{عابر پیاده ('âber-e piyâde)}]{\fa{عابر} \fa{پیاده} ('âber-e piyâde) — a pedestrian}

\label{lesson:FA-C428-aber-e-piyade}



\begin{quote}
Before the new one: say the Persian for a tow truck, then the Persian for to tow.
\end{quote}

\subsection*{You'll want to know: \fa{عابر} \fa{پیاده}}
\textbf{\fa{عابر} \fa{پیاده}} (\emph{'âber-e piyâde}) — \textquotedblleft{}a pedestrian\textquotedblright{}.

\begin{quote}
\fa{راننده} \fa{برای} \fa{عابر} \fa{پیاده} \fa{ایستاد}.
\end{quote}

\begin{quote}
\emph{Rânande barâ-ye 'âber-e piyâde istâd.}
\end{quote}

\begin{quote}
The driver stopped for the pedestrian.
\end{quote}

\emph{'âber} 'passer-by' + \emph{piyâde} 'on foot'. Plural \emph{'âberân-e piyâde}.

\begin{cousinweb}
\emph{'âber-e piyâde} is built from \emph{'âber} 'one who passes', a loan usually traced to Arabic \emph{ʿābir} 'crossing, passing', and native \emph{piyâde} 'on foot', which EtymDB traces to Middle Persian \emph{payādag}.
\end{cousinweb}

\begin{grammarlens}[title={What you've built}]
One more word for telling, explaining and describing at B1.
\end{grammarlens}

\subsection*{Your turn}
\begin{itemize}
\raggedright
  \item \emph{Say it:} \emph{'âber-e piyâde}
  \item \emph{Say it:} \emph{'âber-e piyâde}, once more
  \item \emph{Say it:} a sentence of your own with \emph{'âber-e piyâde}
  \item \emph{Recall:} say the Persian for a tow truck, then the Persian for to tow, then say \emph{'âber-e piyâde} again
\end{itemize}

\begin{tcolorbox}[breakable,colback=teal!4,colframe=teal!35!black,title={Before you move on}]
What does \fa{عابر} \fa{پیاده} mean? (\textquotedblleft{}A pedestrian\textquotedblright{}.) Say it once more.
\end{tcolorbox}
\section[\texorpdfstring{\fa{موتورسیکلت} (motorsiklet)}{موتورسیکلت (motorsiklet)}]{\fa{موتورسیکلت} (motorsiklet) — a motorcycle, a motorbike}

\label{lesson:FA-C428-motorsiklet}



\begin{quote}
Before the new one: say the Persian for to tow, then the Persian for a pedestrian.
\end{quote}

\subsection*{You'll want to know: \fa{موتورسیکلت}}
\textbf{\fa{موتورسیکلت}} (\emph{motorsiklet}) — \textquotedblleft{}a motorcycle, a motorbike\textquotedblright{}.

\begin{quote}
\fa{پسر} \fa{همسایه} \fa{با} \fa{موتورسیکلت} \fa{به} \fa{مدرسه} \fa{میرود}.
\end{quote}

\begin{quote}
\emph{Pesar-e hamsâye bâ motorsiklet be madrese miravad.}
\end{quote}

\begin{quote}
The neighbour's son goes to school on a motorbike.
\end{quote}

In speech usually shortened to \emph{motor}.

\begin{cousinweb}
\emph{motorsiklet} was borrowed from a European language and is usually traced to French \emph{motocyclette} or English \emph{motorcycle}. EtymDB does not list the Persian word.
\end{cousinweb}

\begin{grammarlens}[title={What you've built}]
One more word for telling, explaining and describing at B1.
\end{grammarlens}

\subsection*{Your turn}
\begin{itemize}
\raggedright
  \item \emph{Say it:} \emph{motorsiklet}
  \item \emph{Say it:} \emph{motorsiklet}, once more
  \item \emph{Say it:} a sentence of your own with \emph{motorsiklet}
  \item \emph{Recall:} say the Persian for to tow, then the Persian for a pedestrian, then say \emph{motorsiklet} again
\end{itemize}

\begin{tcolorbox}[breakable,colback=teal!4,colframe=teal!35!black,title={Before you move on}]
What does \fa{موتورسیکلت} mean? (\textquotedblleft{}A motorcycle, a motorbike\textquotedblright{}.) Say it once more.
\end{tcolorbox}
\section[\texorpdfstring{\fa{کامیون} (kâmiyun)}{کامیون (kâmiyun)}]{\fa{کامیون} (kâmiyun) — a lorry (US: a truck)}

\label{lesson:FA-C428-kamiyun}



\begin{quote}
Before the new one: say the Persian for a pedestrian, then the Persian for a motorcycle.
\end{quote}

\subsection*{You'll want to know: \fa{کامیون}}
\textbf{\fa{کامیون}} (\emph{kâmiyun}) — \textquotedblleft{}a lorry (US: a truck)\textquotedblright{}.

\begin{quote}
\fa{کامیونها} \fa{شب} \fa{در} \fa{جاده} \fa{حرکت} \fa{میکنند}.
\end{quote}

\begin{quote}
\emph{Kâmiyun-hâ shab dar jâdde harekat mikonand.}
\end{quote}

\begin{quote}
The lorries travel on the road at night.
\end{quote}

A lorry driver is a \emph{kâmiyundâr} or \emph{rânande-ye kâmiyun}.

\begin{cousinweb}
\emph{kâmiyun} was borrowed from French \emph{camion} 'truck, lorry' (EtymDB).
\end{cousinweb}

\begin{grammarlens}[title={What you've built}]
One more word for telling, explaining and describing at B1.
\end{grammarlens}

\subsection*{Your turn}
\begin{itemize}
\raggedright
  \item \emph{Say it:} \emph{kâmiyun}
  \item \emph{Say it:} \emph{kâmiyun}, once more
  \item \emph{Say it:} a sentence of your own with \emph{kâmiyun}
  \item \emph{Recall:} say the Persian for a pedestrian, then the Persian for a motorcycle, then say \emph{kâmiyun} again
\end{itemize}

\begin{tcolorbox}[breakable,colback=teal!4,colframe=teal!35!black,title={Before you move on}]
What does \fa{کامیون} mean? (\textquotedblleft{}A lorry (US: a truck)\textquotedblright{}.) Say it once more.
\end{tcolorbox}
\section[\texorpdfstring{\fa{وانت} (vânet)}{وانت (vânet)}]{\fa{وانت} (vânet) — a pickup (truck), a small van}

\label{lesson:FA-C428-vanet}



\begin{quote}
Before the new one: say the Persian for a motorcycle, then the Persian for a lorry.
\end{quote}

\subsection*{You'll want to know: \fa{وانت}}
\textbf{\fa{وانت}} (\emph{vânet}) — \textquotedblleft{}a pickup (truck), a small van\textquotedblright{}.

\begin{quote}
\fa{میوه} \fa{فروش} \fa{میوه} \fa{را} \fa{با} \fa{وانت} \fa{آورد}.
\end{quote}

\begin{quote}
\emph{Mive-forush mive râ bâ vânet âvard.}
\end{quote}

\begin{quote}
The fruit seller brought the fruit in a pickup.
\end{quote}

Vans like this carry goods, furniture and fruit all over Iranian towns.

\begin{cousinweb}
\emph{vânet} is a loan word whose exact source is not settled: it is usually linked with English \emph{van}, a small goods vehicle. EtymDB does not list it.
\end{cousinweb}

\begin{grammarlens}[title={What you've built}]
One more word for telling, explaining and describing at B1.
\end{grammarlens}

\subsection*{Your turn}
\begin{itemize}
\raggedright
  \item \emph{Say it:} \emph{vânet}
  \item \emph{Say it:} \emph{vânet}, once more
  \item \emph{Say it:} a sentence of your own with \emph{vânet}
  \item \emph{Recall:} say the Persian for a motorcycle, then the Persian for a lorry, then say \emph{vânet} again
\end{itemize}

\begin{tcolorbox}[breakable,colback=teal!4,colframe=teal!35!black,title={Before you move on}]
What does \fa{وانت} mean? (\textquotedblleft{}A pickup (truck), a small van\textquotedblright{}.) Say it once more.
\end{tcolorbox}
\section[\texorpdfstring{\fa{زیر} \fa{گرفتن} (zir gereftan)}{زیر گرفتن (zir gereftan)}]{\fa{زیر} \fa{گرفتن} (zir gereftan) — to run over, to knock down (with a vehicle)}

\label{lesson:FA-C428-zir-gereftan}



\begin{quote}
Before the new one: say the Persian for a lorry, then the Persian for a pickup.
\end{quote}

\subsection*{You'll want to know: \fa{زیر} \fa{گرفتن}}
\textbf{\fa{زیر} \fa{گرفتن}} (\emph{zir gereftan}) — \textquotedblleft{}to run over, to knock down (with a vehicle)\textquotedblright{}.

\begin{quote}
\fa{راننده} \fa{نزدیک} \fa{بود} \fa{یک} \fa{گربه} \fa{را} \fa{زیر} \fa{بگیرد}.
\end{quote}

\begin{quote}
\emph{Rânande nazdik bud yek gorbe râ zir begirad.}
\end{quote}

\begin{quote}
The driver nearly ran over a cat.
\end{quote}

\emph{zir} 'under' + \emph{gereftan} 'to take'. Present stem \emph{gir-}; \emph{nazdik bud} + subjunctive — nearly.

\begin{cousinweb}
\emph{zir gereftan} is built from native \emph{zir} 'under', which EtymDB lists, and \emph{gereftan} 'to take', usually traced to Old Persian \emph{grab-} 'to seize'.
\end{cousinweb}

\begin{grammarlens}[title={What you've built}]
One more word for telling, explaining and describing at B1.
\end{grammarlens}

\subsection*{Your turn}
\begin{itemize}
\raggedright
  \item \emph{Say it:} \emph{zir gereftan}
  \item \emph{Say it:} \emph{zir gereftan}, once more
  \item \emph{Say it:} a sentence of your own with \emph{zir gereftan}
  \item \emph{Recall:} say the Persian for a lorry, then the Persian for a pickup, then say \emph{zir gereftan} again
\end{itemize}

\begin{tcolorbox}[breakable,colback=teal!4,colframe=teal!35!black,title={Before you move on}]
What does \fa{زیر} \fa{گرفتن} mean? (\textquotedblleft{}To run over, to knock down (with a vehicle)\textquotedblright{}.) Say it once more.
\end{tcolorbox}
